\documentclass[11pt]{article}
\usepackage[utf8]{inputenc}
\usepackage[T1]{fontenc}
\usepackage{lmodern}
\usepackage{textcomp}
\usepackage{amsmath, amssymb, amsthm}
\usepackage{graphicx}
\usepackage{listings}
\usepackage[hidelinks]{hyperref}

\theoremstyle{plain}
\newtheorem{theorem}{Teorema}[section]
\newtheorem{lemma}[theorem]{Lema}
\newtheorem{proposition}[theorem]{Proposición}
\newtheorem{corollary}[theorem]{Corolario}

\theoremstyle{definition}
\newtheorem{definition}[theorem]{Definición}


\title{c4}
\date{}
\begin{document}
\maketitle

\section{Chapter 4}

\section{Mathematical Sentences}

Consider the following theorem of analysis:

Let f : R $\to$R be a differentiable function. Then f is continuous. (4.1)

This statement comprises two sentences. The first sentence does not state a fact: it's an assumption. We aren't told which function f is, so there is no question of this statement about f being true or false. But still we use it as a basis for the rest of the argument. The second sentence does just that: it's a deduction of a new fact from the assumption. The symbol f is an internal variable of the statement, and we can make it disappear:

Every differentiable real function is continuous.

This statement is equivalent to (4.1).

In this chapter we begin to analyse mathematical sentences, and for this purpose we need some elements of logic. The chief attribute of a mathematical statement is its truth or falsehood. Accordingly, we define the two logical (or boolean) constants True and False, abbreviated t, f or 1, 0, respectively. The simplest sentences are the relational expressions, formed by combining pairs of mathematical objects via relational operators. The value of a relational expression is a logical constant. To form more complex sentences we combine relational expressions using logical operators, much like combining numbers using addition and multiplication. The formulation of high-level statements such as (4.1) requires logical functions---called predicates---and quantifiers, which resemble integrals.

Theseconstructshaveaccompanyingsymbols,andthesymbolsoflogicarecryptic and alluring. While they help us understand the structure of sentences, they don't necessarily help us write good sentences. Excessive use of logical symbols clutters the exposition, obscuring meaning. A recurrent theme of this chapter is that, quite often, concepts of great logical depth are better expressed with words.

© Springer-Verlag London 2014 F. Vivaldi, Mathematical Writing, Springer Undergraduate Mathematics Series, DOI 10.1007/978-1-4471-6527-9\_4

51

52 4 Mathematical Sentences

\section{4.1 Relational Operators}

We begin with a simple symbolic sentence:

0 < 1. (4.2)

This is a relational expression. The relational operator `<' converts two numbers into the logical constant True. Any significant mathematical sentence that evaluates to True is called a theorem, and expression (4.2) is one of the first theorems of analysis; it underpins all inequalities among real numbers.

Relational operators comprise very familiar objects:

$\neq$ = < $\leq$ > $\geq$. (4.3)

These operators are binary, they act on two operands. The first two act on elements of any set; the others act on real numbers (more generally, on elements of an ordered set---see Sect.4.6).

Interesting things can be said with simple relational expressions:

< 3 $\times$ 10$-$7.

93 + 103 = 13 + 123, $\pi$ $-$355

(4.4)

113

These expressions also feature prominently in programming languages:

if x<1 then

x := -x + 1 else

x := x - 1 fi:

Here the logical value of the expression `x < 1' determines which assignment statement is executed.

At the most basic level we have the relational operators associated with sets, namely the membership and inclusion operators (Sect.2.1) and their negation:

$\notin$ $\in$  $\subset$  $\supset$. (4.5)

For example, the relational expression

$\surd$

2 $\in$Q (4.6)

is True, that is, the square root of 2 is not a rational number. We shall prove this fact in Sect.7.1.

There are countless relational operators in mathematics: the divisibility operator `|' and the congruence operator `$\equiv$' in arithmetic (Sect.2.1.3), the isomorphism operator `$\sim$=' in algebra, the orthogonality operator `$\perp$' in geometry, etc. By

4.1 Relational Operators 53

contrast,thesymbol`$\approx$',usedinmathematicalphysics,doesnotrepresentarelational operator because expressions such as $\pi$ $\approx$3.14 cannot be assigned unequivocally a value True or False.

\section{4.2 Logical Operators}

The following sentence is True:

Theinteger 257885161 $-$1 is prime. (4.7)

The proof requires a sophisticated theory and a lot of computer time.1 However, it is not difficult to see that this expression is decomposable into finitely many relational expressions. Indeed the primality of p = 2 57885161 $-$1 could in principle be established by verifying that p is not divisible by any prime up to $\surd$p. (The

locution `in principle' is cheeky; no computer present or future will ever complete such a brute-force computation within the lifetime of our Universe.)

This example motivates the introduction of the logical (or boolean) operators not and and (to be defined below), with which we rewrite (4.7) as follows:

not (2|p) and not (3|p) and not (5|p) and \textperiodcentered{} \textperiodcentered{} \textperiodcentered{} .

We have produced a complex symbolic sentence by combining relational expressions with logical operators. This process is analogous to the construction of complex arithmetical expressions from arithmetical constants (numbers) and operators (+, $-$, etc.).

The operator not (represented by the symbol $\neg$) is unary---it takes one boolean operand and produces a boolean value. In this respect not resembles the unary arithmetical operator `$-$', which changes the sign of a number. The operator and (represented by the symbol $\land$) is binary: think of it as a kind of multiplication.

The following truth tables define the operators not and and by specifying their action on all possible choices of operands:

P Q P $\land$Q T T T T F F F T F F F F

P $\neg$P T F F T

(4.8)

The expression $\neg$P is called the negation of P. Negating relational expressions is straightforward, as we already have all the relevant symbols---see (4.3) and (4.5).

1 This expression refers to the largest known prime (as of January 2014) with 17425170 digits: see [6].

54 4 Mathematical Sentences

Thus

$\neg$(x < y) = (x $\geq$y), $\neg$(x $\in$A) = (x $\in$A), $\neg$(A $\subset$B) = A $\subset$B.

The expression P $\land$Q is called a conjunction or compound expression. In a compound expression the operator and may appear implicitly. For instance, the expression 0 < x < 1 is compound: (0 < x) $\land$(x < 1).

Next we define the operator or (represented by the symbol $\lor$) and the implication operator (represented by the symbol $\Rightarrow$), as follows:

P Q P $\lor$Q T T T T F T F T T F F F

P Q P $\Rightarrow$Q T T T T F F F T T F F T

(4.9)

These operators could also be expressed in terms of $\neg$ and $\land$, see Exercise 4.10. The operator or is inclusive, namely t $\lor$t = t, which is not the meaning usually attributed to the conjunction `or' in English usage.2

The expression

P $\Rightarrow$Q (4.10)

is called an implication, and it reads

P implies Q Q follows from P if P, then Q

P only if Q P is sufficient for Q.

The expression P is the hypothesis (or antecedent) and Q is the conclusion (or consequent). Many mathematical statements have this form, for example theorem (4.1).

According to the truth table (4.9), if P is false, then the expression P $\Rightarrow$Q is true for any Q. This convention, which is perhaps unexpected, is in fact in agreement with the common usage of an implication:

If 4 is a prime number, then I am a Martian. If I win the lottery, then I'll buy you a Ferrari.

The meaning of these sentences is clear. In particular, if I don't win the lottery, then I am free do what I want about the Ferrari.

From (4.9) we see that the expressions P $\Rightarrow$Q and Q $\Rightarrow$P are different. Accordingly, we introduce the operator $\Leftarrow$as follows:

P $\Leftarrow$Q def = Q $\Rightarrow$P.

2 There is also an exclusive version of OR, called XOR (), for which T T = F.

4.2 Logical Operators 55

The expression on the left is called the converse of the implication (4.10). It reads

P is implied by Q P follows from Q

P if Q P is necessary for Q.

The value of an implication and its converse are unrelated. In the following example the former is false and the latter is true:

(x2 = 25) $\Rightarrow$(x = $-$5) (x2 = 25) $\Leftarrow$(x = $-$5). (4.11)

Inappropriate reversal of an implication is a common mistake in proofs---see Sect.7.7.2.

Formulae (4.11) combine relational and boolean operators, and the parentheses suggest the order in which these operators are evaluated. In this case the parentheses are redundant, since relational operators must necessarily take precedence over boolean operators. Thus the expression x2 = 25 $\Rightarrow$x = $-$5 could not be interpreted as x2 = (25 $\Rightarrow$x) = $-$5. However, when two or more boolean operators are present, parentheses may be inserted to change the meaning of an expression. For instance, the expressions

(f $\land$t) $\Rightarrow$t f $\land$(t $\Rightarrow$t)

evaluate to True and False respectively. As with arithmetical operators, there is an agreed order with which boolean operators are evaluated: first $\neg$, then $\land$, then $\lor$, then $\Rightarrow$. Thus the expression P$\lor$$\neg$Q $\Rightarrow$R$\land$S is evaluated as (P$\lor$($\neg$Q)) $\Rightarrow$(R$\land$S); we note that the redundant parentheses add clarity to the sentence.

The equivalence operator $\Leftrightarrow$is the conjunction of the direct and converse implications:

P $\Leftrightarrow$Q def = (P $\Rightarrow$Q) $\land$(P $\Leftarrow$Q) (4.12)

and it corresponds to the following truth table:

P Q P $\Leftrightarrow$Q T T T T F F F T F F F T

The expression P $\Leftrightarrow$Q is read aloud as

P implies and is implied by Q P is equivalent to Q

P if and only if Q P is necessary and sufficient for Q

The awkward expression `if and only if' is quite common; the abbreviated form `iff' is found in formulae as an alternative to `$\Leftrightarrow$'. So the expressions

56 4 Mathematical Sentences

A = B $\Leftrightarrow$A $\Delta$ B = $\emptyset$ A = B iff A $\Delta$ B = $\emptyset$

have the same meaning; we turn symbols into words:

Two sets are equal if and only if their symmetric difference is empty.

If---as in this example---the value of P $\Leftrightarrow$Q is True, then the statements P and Q are logically equivalent, meaning that one is is just a rewording of the other.

The contrapositive of the implication (4.10) is the implication

$\neg$P $\Leftarrow$$\neg$Q. (4.13)

Thecontrapositiveis constructedbyreversingtheoperator andnegatingtheoperands. Great care must be exercised in distinguishing between direct, converse, and contrapositive implications.

Direct : If x is a multiple of 4, then x is even. (True) Converse : If x is even, then x is a multiple of 4. (False) Contrapositive : If x is odd then x is not a multiple of 4. (True)

While the value of an implication and of its converse are unrelated, every implication is equivalent to its contrapositive. This is identity (iv) of the following theorem.

\begin{theorem}
Theorem 4.1 For all P, Q $\in$\{t, f\}, the following holds:
\end{theorem}

(i) $\neg$(P $\lor$Q) $\Leftrightarrow$($\neg$P $\land$$\neg$Q) (ii) $\neg$(P $\land$Q) $\Leftrightarrow$($\neg$P $\lor$$\neg$Q) (iii) (P $\Rightarrow$Q) $\Leftrightarrow$($\neg$P $\lor$Q) (iv) (P $\Rightarrow$Q) $\Leftrightarrow$($\neg$P $\Leftarrow$$\neg$Q).

\begin{proof}
Proof. We prove (iii), leaving the other cases as an exercise. We will show that the two sides of the equivalence (iii) have the same truth table. The truth table of the left-hand side is given in (4.9); that of the right-hand is computed as follows:
\end{proof}

P Q $\neg$P $\neg$P $\lor$Q P $\Rightarrow$Q T T F T T T F F F F F T T T T F F T T T

We see that the two truth tables are equal. 

The statements (i) and (ii) are known as De Morgan's laws.3 Using Theorem 4.1, one can express the operators $\lor$, $\Rightarrow$, $\Leftrightarrow$in terms of $\neg$ and $\land$(Exercise 4.10).

All items in the theorem have the form A(P, Q) $\Leftrightarrow$B(P, Q), where the statements A and B have the same value for any choice of P and Q and hence are equivalent. A statement of this kind is called a tautology. In logic a distinction is made between a

3 Augustus De Morgan (British: 1806--1871).

4.2 Logical Operators 57

tautology and a syntactically correct---but not necessarily true---double implication, by introducing the symbol $\leftrightarrow$for the latter circumstance. So one would write (P $\lor$Q) $\leftrightarrow$(P $\land$Q) rather than (P $\lor$Q) $\Leftrightarrow$(P $\land$Q). A similar distinction exists between $\Rightarrow$and $\to$. For our purpose these additional symbols are unnecessary.

\section{4.3 Predicates}

The sentence

The integer x is divisible by 7

contains an unspecified symbol, and thus it assumes a logical value only if the symbol is assigned an integer value. We are prompted to introduce the following function:

P : Z $\to$\{t, f\} x $\to$7 | x. (4.14)

We see that P(x) is the symbolic translation of our sentence, and that P(22) = f, P($-$91) = t.

Using a function to encode a sentence is a simple idea with far-reaching consequences. Let us thus define a predicate (or boolean function) to be any function that assumes boolean values. If the domain of a predicate P is X, then we speak of a predicate on (or over) X. So (4.14) is a predicate on the integers.

Let P : X $\to$\{t, f\} be a predicate on X. There is a distinguished subset AP of X, determined by P via the following Zermelo definition:

AP := \{x $\in$X : P(x)\}. (4.15)

With reference to Eq. (2.9) and the discussion that follows, we see that the expression P(x), which means `x has property P', is the definiens of the set AP. So the Zermelo definition (4.15) may be expressed in the language of functions using an inverse image:

\{x $\in$X : P(x)\} def = P$-$1(\{t\})

[cf. Definition (2.20)].

Conversely, let X be a set and let A be a subset of X. The predicate

PA : X $\to$\{t, f\} x $\to$x $\in$A (4.16)

is called the characteristic function of A (in X). Explicitly, P(x) is equal to t if x belongs to A and to f otherwise. For example, the function (4.14) is the characteristic function of the set 7Z of integer multiples of 7.

So to every predicate on a set X we associate a unique subset of X, and vice-versa. We say that there is a bi-unique correspondence between these two classes of objects, established by expressions (4.15) and (4.16).

58 4 Mathematical Sentences

Setsaremanipulatedusingsetoperators.Predicatesaremanipulatedusingboolean operators in a natural way by defining, say,

(P $\land$Q)(x) def = P(x) $\land$Q(x)

and similarly for all other operators. The following theorem establishes correspondences between these two classes of objects.

\begin{theorem}
Theorem 4.2 Let X be a set, let A, B $\subseteq$X, and let PA, PB be the corresponding characteristic functions. The following holds (the prime denotes taking the complement):
\end{theorem}

(i) $\neg$ PA = PA $'$ (ii) PA $\land$PB = PA $\cap$B (iii) PA $\lor$PB = PA $\cup$B (iv) PA$\Rightarrow$PB = P(A$\setminus$B)$'$ (v) PA$\Leftrightarrow$PB = P(A$\cap$B) $\cup$(A $\cup$B)$'$.

(This baroque display of symbols is unappealing. In Sect.6.6 we'll make it more digestible by writing a short essay about it without symbols.)

\begin{proof}
Proof. To prove (i) we note that the function x $\to$$\neg$PA(x) evaluates to True if x $\in$A and to False otherwise. However, from the definition of the complement of a set, we have x $\in$A $\Leftrightarrow$x $\in$A$'$.
\end{proof}

Next we prove (iv). We'll prove instead

$\neg$(PA $\Rightarrow$PB) = P(A$\setminus$B)

which, together with (i), gives us (iv). Let PA := (x $\in$A) and let PB := (x $\in$ B). From the truth table (4.9) of the operator $\Rightarrow$, we find that $\neg$(PA $\Rightarrow$PB)(x) is True precisely when PA(x) is True and PB(x) is False. This means that

(x $\in$A) $\land$(x $\in$B),

but this is just the definition of the characteristic function of the set A$\setminus$B, as desired.

The proof of (ii), (iii), (v) is left as an exercise. 

We illustrate the significance of this theorem with examples.

Example.Letaandbberealnumbers.Thepredicatesx $\to$(x $\geq$a)andx $\to$(x < b) (over R) are the characteristic functions of two rays, the latter without end-point. According to Theorem 4.2 part (ii), the predicate x $\to$((x $\geq$a) $\land$(x < b)) is the characteristic function of the intersection of these rays. Depending on whether a < b, or a $\geq$b, this intersection is the half-open interval [a, b) or the empty set.

Example. Let X, Y be sets, and let f , g : X $\to$Y be functions. An equation (on X) is a predicate of the type

4.3 Predicates 59

P : X $\to$\{t, f\} x $\to$( f (x) = g(x)).

In the language of predicates, an equation is the characteristic function of its own solution set \{x $\in$X : P(x)\} (Sect.3.3). The system of two equations

f1(x) = g1(x) f2(x) = g2(x)

defined over the same set corresponds to the predicate x $\to$P1(x) $\land$P2(x). From Theorem 4.2, part (ii), it follows that the solution set of a system of two equations is the intersection of the solution sets of the individual equations. The same holds for any number of equations---see (3.19).

Example. The standard form of the sigma-notation (see Sect.3.2) has the following syntax:

ak (4.17)

P(k)

where P is any predicate over Z. The range of summation consists of those values of k for which P(k) is True.

Example. If A $\subset$B, then A$\setminus$B is empty, and from part (iv) of the theorem we obtain

(PA $\Rightarrow$PB) = P$\emptyset$$'$ = PX.

The subset relation has been translated into an identity of functions. For example, the functional identity

(P4Z $\Rightarrow$P2Z) = PZ

says that every multiple of four is even.

\section{4.4 Quantifiers}

Interesting mathematical statements---such as (4.1)---refer to families of objects rather than to individual objects. The formulation of general statements requires two special symbols, the universal quantifier $\forall$and the existential quantifier $\exists$. These symbols translate into words in several equivalent ways:

$\forall$: for all given any for any choice of

$\exists$: for some there exists we can find

Expressions with quantifiers have the following syntax:

$\forall$x $\in$X, P(x) For all x in X, P(x).

(4.18) $\exists$x $\in$X, P(x) There exists x in X, such that P(x).

60 4 Mathematical Sentences

where X is a set and P is a predicate over X. The quantifier is followed by the symbol being quantified, the membership operator, a set, and a predicate. The acronym `s.t.' (for `such that') is sometimes inserted in the second symbolic expression (4.18) before the predicate:

$\exists$x $\in$X s.t. P(x).

Without a predicate the expression $\forall$x $\in$X is incomplete and meaningless. The expression $\exists$x $\in$X is meaningful but not used; one instead writes X = $\emptyset$.

The meaning of the expressions (4.18) should be clear; in particular, each expression has a logical value. For example, the following sentences are True:

$\forall$x $\in$R, x2 $\geq$0 The square of every real number is non-negative. $\exists$x $\in$Q, 0 < x < 1 There is a rational number in the open unit interval.

In translating the symbolic expressions we haven't merely converted symbols into words following (4.18). We have also synthesised meaning, and in doing so the quantified variable x has become silent---no explicit reference to it remains in the sentence. This is an intrinsic phenomenon, as we shall see.

If the ambient set is clear from the context, then reference to it may be omitted. This is invariably the case if the quantifier appears in conjunction with an inequality that restricts the variable's range. Of the following equivalent expressions

$\forall$n > 3, n! > 2n $\forall$n $\in$N$\setminus$\{1, 2, 3\}, n! > 2n $\forall$n $\in$N, n > 3 $\Rightarrow$n! > 2n

the first one is preferable. The presence of the factorial makes it clear that we are dealing with natural numbers.

Example. Using quantifiers, it is possible to define relational operators in terms of other operators, thereby reducing the number of primitive operators. The inclusion operator can be defined in terms of the membership operator:

A $\subset$B $\Leftrightarrow$ def $\forall$x $\in$A, x $\in$B. (4.19)

The divisibility operator is defined in terms of multiplication:

m|n $\Leftrightarrow$ def $\exists$k $\in$Z, mk = n. (4.20)

Example. The symbolic sentences (4.18) may be expressed in the language of functions:

$\forall$x $\in$X, P(x) $\Leftrightarrow$P(X) = \{t\}

(4.21)

$\exists$x $\in$X, P(x) $\Leftrightarrow$P(X) = \{f\}.

These equivalences may give the impression that we have managed to bypass quantifiers altogether. This is not the case. The standard definition of image of a set under

4.4 Quantifiers 61

a function (2.18) has a hidden quantifier in it, which is revealed when the definition is converted to Zermelo form. Thus if f : A $\to$B is a function, then

f (A) = \{f (x) : x $\in$A\} = \{y $\in$B : $\exists$x $\in$A, y = f (x)\}.

Of the two definitions, the first one is more intuitive.

Several quantifiers may appear within the same sentence. For illustration, let us consider the following two-variable predicate:

L (x, y) := (`x loves y').

This function is defined over the cartesian product G$\times$G, where G is a set of people. So the lovers and the loved ones belong to the same set.

Now choose g $\in$G, say, g = George. Then `everybody loves George' is a boolean expression, which is spelled out as

For all elements x of G, the value of L (x, g) is True.

This expression translates into symbols as follows:

$\forall$x $\in$G, L (x, g).

This sentence may be true or false (depending on how charming George is). The followingsentenceswithoneortwoquantifiersexemplifyseveralpossibleconstructs.

Words Symbols

1 Everybody loves George $\forall$x $\in$G, L (x, g)

2 Somebody loves George $\exists$x $\in$G, L (x, g)

3 Everybody loves himself $\forall$x $\in$G, L (x, x)

4 George is in love $\exists$x $\in$G$\setminus$\{g\}, L (g, x)

5 George is selfish L (g, g) $\land$($\forall$x $\in$G$\setminus$\{g\}, $\neg$L (g, x))

6 Everybody is in love $\forall$x $\in$G, $\exists$y $\in$G, L (x, y)

7 Somebody is in love $\exists$x $\in$G, $\exists$y $\in$G, L (x, y)

8 Somebody loves everybody $\exists$x $\in$G, $\forall$y $\in$G, L (x, y)

9 Somebody is loved by everybody $\exists$y $\in$G, $\forall$x $\in$G, L (x, y)

In the sentences 4 and 5, the indeterminate x belongs to the set G$\setminus$\{g\}, to ensure that x is distinct from g. It is possible to transfer this constraint to the predicate, e.g., in item 4:

$\exists$x $\in$G, (x = g) $\land$L (g, x).

62 4 Mathematical Sentences

Two-quantifier sentences have implied parentheses, which establish the order of evaluation of the sub-expressions. Thus expression 6 is written in full as

$\forall$x $\in$G, ($\exists$y $\in$G, L (x, y)).

According to (4.18), the quantity in parentheses must be equal to Q(x) for some predicate Q over G; indeed we see that Q(x) =`x loves somebody'. From the sentence's symbolic structure and meaning, it is clear that the choice of y depends on x, in general. Therefore, exchanging the order of the quantifiers alters the meaning of the expression, as confirmed by comparing sentences 6 and 9.

Quantifiers are operators which act on boolean functions, thereby producing new functions. Their effect is to reduce the number of variables by one, a process that calls to mind definite integration. We should think of the two expressions

$\forall$x $\in$X

dx

X

as being structurally similar. The operator on the left acts on predicates over the set X, namely on functions P : X $\to$\{t, f\}. The operator on the right acts, say, on functions F : X $\to$R, which are integrable over a subset X of the real line. (To fix ideas, assume that X = [a, b] is an interval.) Inserting the appropriate functions in each expression

dx F(x)

$\forall$x $\in$X, P(x)

X

we obtain a boolean constant (True or False) on the left, and a numerical constant (a number) on the right. Now suppose that both P and F are functions of two variables x and y, defined over the cartesian product X $\times$ Y of two sets. Then each expression

dx F(x, y)

$\forall$x $\in$X, P(x, y)

X

produces a function of y, the variable that is not quantified in one case, and not integrated over in the other. If we quantify/integrate with respect to both variables,

dy F(x, y),

$\forall$x $\in$X, $\exists$y $\in$Y, P(x, y)

dx

X

Y

once again we obtain constants.

Returning to the examples above, we note that L (x, g) is a predicate in one variable (g is fixed), and so is L (x, x). Hence $\exists$x, L (x, g) and $\forall$x, L (x, x) are constants. On the other hand, any expression in which the number of quantifiers is smaller than the number of arguments in the predicate is a predicate with fewer arguments, as the following examples illustrate.

4.4 Quantifiers 63

Words Symbols

x is in love $\exists$y $\in$G$\setminus$\{x\}, L (x, y)

x is loved $\exists$y $\in$G$\setminus$\{x\}, L (y, x)

x is a hippy $\forall$y $\in$G, L (x, y)

x is selfish L (x, x) $\land$($\forall$y $\in$G$\setminus$\{x\}, $\neg$L (x, y))

x is a lover of George x = g $\land$L (x, g) $\land$L (g, x)

Now, each predicate is the characteristic function of a subset of G, which we construct with a Zermelo definition.

Words Symbols

the people in love \{x $\in$G : $\exists$y $\in$G$\setminus$\{x\}, L (x, y)\}

the loved ones \{x $\in$G : $\exists$y $\in$G$\setminus$\{x\}, L (y, x)\}

the hippies \{x $\in$G : $\forall$y $\in$G, L (x, y)\}

the selfish people \{x $\in$G : L (x, x) $\land$($\forall$y $\in$G$\setminus$\{x\}, $\neg$L (x, y))\}

George$'$s lovers \{x $\in$G : x = g $\land$L (x, g) $\land$L (g, x)\}

The sentence `everybody loves somebody' is of the form

$\forall$x $\in$X, $\exists$y $\in$Y, P(x, y) (4.22)

where X and Y are sets and P is a predicate over X $\times$ Y. Letting X = Y = N and P(x, y) = (x < y), we obtain a new statement which has the same formal structure:

Given any integer, one can find a larger integer.

or better,

There is no greatest integer.

It's useful to think of the interplay between universal and existential quantifiers as a representation of an adversarial system---like a court of law---based on the following rules of engagement:

$\forall$ given any (my opponent's move)

$\exists$ I can find (my move).

The quantifiers create a tension between the two contenders, and I should expect to be challenged by my opponent. Accordingly, we rewrite our statement more emphatically:

Given any integer, no matter how large, one can always find a larger integer.

The expressions `no matter how large' and `always' are inessential to the claim, but they expose the dynamics of the process. This aspect of predicate calculus will be considered again in Chap.7, when we deal with proofs.

64 4 Mathematical Sentences

Example. The Archimedean property of the real numbers states that

The integer multiples of any positive quantity can be made arbitrarily large.

The symbolic version of this statement requires three quantifiers.

$\forall$x $\in$R+, $\forall$y $\in$R, $\exists$n $\in$N, nx > y.

In this expression x is the positive quantity in question, y is the (large) quantity we want to exceed, and n is the multiple of x needed to achieve this. Note that n = n(x, y). Clearly, the Archimedean property is better expressed with words than with symbols, and the large number of quantifiers hidden within this sentence may be taken as an indication of the concept's logical depth.

\section{4.4.1 Quantifiers and Functions}

Invariably, statements about functions will refer to the elements of the domain or co-domain, and quantifiers may be made to act on these elements. For instance, a function is injective if it maps distinct points to distinct points (Sect.2.2). This concise statement unfolds as follows:

Given any two points in the domain of the function, if they are distinct, then so are their images under the function.

This sentence has a universal quantifier (`given any') and an implication operator (`if . . . then'). Let f : A $\to$B. A literal translation into symbols is

$\forall$x, y $\in$A, (x = y) $\Rightarrow$(f (x) = f (y))

where $\forall$x, y $\in$A is a shorthand for $\forall$x $\in$A, $\forall$y $\in$A. Replacing the implication by its contrapositive [Theorem 4.1 (iv)] we obtain a neater expression:

$\forall$x, y $\in$A, (f (x) = f (y)) $\Rightarrow$(x = y).

The original definition (`distinct points have distinct images') restricts the ambient set to pairs of distinct points. This can be done symbolically:

$\forall$x $\in$A, $\forall$y $\in$A$\setminus$\{x\}, f (x) = f (y).

Now the predicate is simpler but the set is more complicated. It is always possible to transform an implication by absorbing the hypothesis into the ambient set (Exercise 4.10.3).

A function f is surjective if image and co-domain coincide: f (A) = B. This high-level expression hides two quantifiers; we first spell it out:

4.4 Quantifiers 65

Given any point in the co-domain, we can find a point in the domain which maps to it.

and then we rewrite it symbolically:

$\forall$y $\in$B, $\exists$x $\in$A, f (x) = y.

This expression is of the form (4.22).

In Sect.3.1 we noted that a sequence of elements of a set A may be interpreted as a function:

a : N $\to$A k $\to$ak.

Characterising sequences is then analogous to characterising functions.

For example, consider the set ZN of all integer sequences (this notation was developed in Sect.3.5.1). To isolate sequences with certain properties---a subset of ZN---we use a Zermelo definition. For example, the set

\{a $\in$ZN : a1 $\in$2Z\}

consists of all integer sequences whose first term is even. By applying the quantifier $\forall$to the subscript k (which is the variable in our functions), we can deal with all terms of the sequences at once. So the set of integer sequences with only even terms is given by

\{a $\in$ZN : $\forall$k $\in$N, ak $\in$2Z\} = (2Z)N. (4.23)

There seems to be something wrong here. The quantifier `integrates out' the indeterminate k: does this mean that the predicate is just a constant? No, because the Zermelo variable is actually a = (ak), an element of the ambient set ZN. If we write (4.23) as \{a $\in$ZN : P(a)\}, then we see the structure of the predicate:

P : ZN $\to$\{t, f\} a $\to$($\forall$k $\in$N, ak $\in$2Z).

Likewise, the set

\{a $\in$ZN : $\exists$k $\in$N, ak $\in$2Z\}

is

The set of integer sequences with an even term.

The above expression reflects the very literal interpretation of what's written, which is what mathematicians---and lawyers---are notorious for. A mathematical geek would be entertained by the fact that the statement `In London there is an underground station', is true; a normal person would instead perceive this as a puzzling under-statement (`Where do they go from there?'). So a characterisation of the type

The set of integer sequences with at least one even term.

66 4 Mathematical Sentences

is preferable to the one given before. The qualifier `at least' is superfluous, but helps the reader note an essential point.

The set

\{a $\in$ZN : $\forall$n $\in$N, n > 2 $\Rightarrow$2|an\}

is described as

The set of integer sequences whose terms, after the second, are even.

No information is given about the first two terms, yet this statement could be misinterpreted as meaning that these terms are odd. A more eloquent description of this set is

The set of integer sequences whose terms are all even, with the possible exception of the first two terms.

These examples show how much can be done to improve the clarity of a definition; this will be our concern in Chap.6.

\section{4.4.2 Existence Statements}

An existence statement is a logical expression with a leading existential quantifier $\exists$. In this section we consider some examples; we shall take a deeper look at the question of existence in Chap.9.

The standard symbolic form of an existence statement is (4.18), but in a mathematical sentence the quantifier is often hidden:

The number cos($\pi$/3) is rational.

To make the quantifier visible, we write

For some rational number r, we have cos($\pi$/3) = r.

or, in symbols,

$\exists$r $\in$Q, r = cos($\pi$/3).

This statement is weaker than the identity cos($\pi$/3) = 1/2, because it does not require us to reveal which rational number our expression is equal to.

Equation (4.20) shows that an existential quantifier is hidden in the definition of the divisibility of integers. Therefore the sentence `n is even', or, more generally, `m divides n' are existence statements.

The sentence `n is not prime' is an existence statement, because it means `there is a proper divisor of n'. We write it symbolically as

$\exists$m $\in$N, m $\in$\{1, n\} $\land$(m|n)

$\Leftrightarrow$$\exists$m $\in$N, m $\in$\{1, n\} $\land$($\exists$k $\in$N, mk = n)

4.4 Quantifiers 67

where we have been careful in excluding trivial divisors.

The following well-known theorem in arithmetic is an existence statement.

Theorem (Lagrange,4 1770). Every natural number is the sum of four squares.

Let us analyse it in detail. First, we consider two instances of the theorem:

5 = 22 + 12 + 02 + 02 7 = 22 + 12 + 12 + 12.

The first identity shows that some integers are the sum of fewer than four squares, while one checks that 7 requires all four squares. We restate the second identity without disclosing the details:

$\exists$a, b, c, d $\in$Z, 7 = a2 + b2 + c2 + d2.

By replacing 5 or 7 with an unspecified natural number n, we obtain a predicate L over N:

$\exists$a, b, c, d $\in$Z, n = a2 + b2 + c2 + d2

L (n) :=

(4.24)

or

n is a sum of four squares.

To state Lagrange's theorem with symbols, we quantify the remaining variable n:

$\forall$n $\in$N, $\exists$a, b, c, d $\in$Z, n = a2 + b2 + c2 + d2. (4.25)

Five quantifiers are necessary to turn the predicate n = a2 +b2 +c2 +d2 (a function of five variables) into a boolean constant. In the original formulation of this theorem, all five variables have fallen silent.

Having buried all meaning inside the predicate (4.25), Lagrange's theorem now disappears into a set identity:

N = \{n $\in$N : L (n)\}.

Any theorem on natural numbers can be put in this form for an appropriate predicate L . (Think about it.)

Negating expressions with a leading universal quantifier always results in existence statements, as we shall see in the next section.

\section{4.5 Negating Logical Expressions}

If L is a logical expression, then its negation is $\neg$L , which is false if L is true and vice-versa. Relational expressions are negated by crossing out the operator with a

4 Joseph-Louis Lagrange, born Giuseppe Lodovico Lagrangia (Italian: 1736--1813).

68 4 Mathematical Sentences

forward slash---see (4.3) and (4.5), Sect.4.1. The only exceptions are the inequalities, whose negations have dedicated symbols. 

\begin{theorem}
Theorem 4.1 provides the negation formulae for the main compound expressions
\end{theorem}

$\neg$(P $\land$Q) $\Leftrightarrow$$\neg$P $\lor$$\neg$Q $\neg$(P $\lor$Q) $\Leftrightarrow$$\neg$P $\land$$\neg$Q (4.26) $\neg$(P $\Rightarrow$Q) $\Leftrightarrow$P $\land$$\neg$Q.

The first two formulae are items (i) and (ii) of the theorem (de Morgan's laws); the third follows immediately from (iii) and (i). Note that the negation of an implication does not have the form of an implication.

If L is an expression with quantifiers and boolean operators, then the expression $\neg$L unfoldsintoanequivalentexpressionwhichweseektodetermine.Thefollowing sentences contain a negation and a universal quantifier:

Not all polynomials have real roots. A square matrix is not necessarily invertible.

We recognise that these are disguised existence statements:

There is a polynomial with no real roots. There is a square matrix which is not invertible.

The following theorem confirms this observation.

\begin{theorem}
Theorem 4.3 Let P be a predicate on a set A, and let
\end{theorem}

L := $\forall$x $\in$A, P(x) M := $\exists$x $\in$A, P(x).

Then

$\neg$L $\Leftrightarrow$$\exists$x $\in$A, $\neg$P(x) $\neg$M $\Leftrightarrow$$\forall$x $\in$A, $\neg$P(x). (4.27)

\begin{proof}
Proof. From (4.21), we have the boolean equivalence
\end{proof}

L $\Leftrightarrow$P(A) = \{t\}

where P(A) is the image of A under P (see Sect.2.2). But then

$\neg$L $\Leftrightarrow$P(A) = \{t\} $\Leftrightarrow$$\neg$P(A) = \{f\}.

Using again (4.21), we obtain

$\neg$L $\Leftrightarrow$$\exists$x $\in$A, $\neg$P(x)

as claimed.

The second formula is proved similarly---see exercises. 

4.5 Negating Logical Expressions 69

It's now easy to deduce the rule for negating expressions with two or more quantifiers. Let P be a predicate on A $\times$ B, and let us consider the expression

L := $\forall$x $\in$A, $\exists$y $\in$B, P(x, y).

Repeated applications of Theorem 4.3 give

$\neg$L $\Leftrightarrow$$\neg$

$\forall$x $\in$A, ($\exists$y $\in$B, P(x, y))

$\Leftrightarrow$$\exists$x $\in$A, $\neg$($\exists$y $\in$B, P(x, y))

$\Leftrightarrow$$\exists$x $\in$A, $\forall$y $\in$B, $\neg$P(x, y).

It should be clear how similar formulae may be derived, having any combination of two quantifiers.

Finally, we consider the case of an arbitrary sequence of quantifiers, e.g.,

L = $\forall$x1 $\in$X1, $\exists$x2 $\in$X2, . . . , $\forall$xn $\in$Xn, P(x1, . . . , xn)

wheretheXi aresets,andP isapredicateonthecartesianproductX1 $\times$ X2 $\times$ \textperiodcentered{} \textperiodcentered{} \textperiodcentered{}$\times$Xn. Parentheses make it clear that this is a nested array of predicates:

L = $\forall$x1 $\in$X1, ( $\exists$x2 $\in$X2, (. . . , ($\forall$xn $\in$Xn, P(x1, . . . , xn)))).

Repeatedly applying Theorem 4.3, from the outside to the inside, we obtain

$\neg$L $\Leftrightarrow$$\exists$x1 $\in$X1, $\forall$x2 $\in$X2, . . . , $\exists$xn $\in$Xn, $\neg$P(x1, . . . , xn).

Namely, the negation of L is obtained by replacing each $\forall$with $\exists$, and vice-versa, without changing their order, and then replacing P with $\neg$P. A formal proof requires the principle of induction (see Chap.8).

Example. The statement

$\forall$n, a, b $\in$Z, (n|ab) $\Rightarrow$(n|a $\lor$n|b) (4.28)

is false (why?). We negate it using Theorem 4.3, obtaining a true statement:

$\exists$n, a, b $\in$Z, (n|ab) $\land$(n | a $\land$n | b).

Example. Lagrange's theorem, expressed symbolically in Eq. (4.25), would be false if four squares were replaced by three squares. This fact is written symbolically as

$\exists$n $\in$N, $\forall$a, b, c $\in$Z, n = a2 + b2 + c2.

In words:

There is a natural number which cannot be written as the sum of three squares.

70 4 Mathematical Sentences

\section{4.6 Relations}

(This section is not essential for the rest of the book: it may be skipped on first reading.)

Relations are special types of logical functions which are ubiquitous in higher mathematics. We introduce some relevant words and symbols.

Let X and Y be sets. A relation R on X $\times$ Y is a predicate over X $\times$ Y. If X = Y, we speak of a relation on X. It is customary to write xRy to mean R(x, y); the expression xRy is called a relational expression, and R is a relational operator. All relational expressions introduced at the beginning of this chapter are of this form. Thus the membership operator $\in$defines a relation on X $\times$ P(X), where X is some ambient set.

A relation on a set X is sometimes defined as a subset of X2, rather than a predicate over X2. In this sense, the set \{(1, 1), (2, 1)\} is a relation on \{1, 2\}. The correspondence between sets and predicates described in Sect.4.3 clarifies the connection between the two constructs.

Among relations on a set X, the equivalence relations hold a special place. They are defined by the following properties:

$\forall$x $\in$X, x R x reflexivity

$\forall$x, y $\in$X, x R y $\Rightarrow$y R x symmetry $\forall$x, y, z $\in$X, (x R y $\land$y R z) $\Rightarrow$x R z transitivity.

The notation x $\sim$y is usually employed to denote equivalence.

Equivalence relations allow us to regroup the elements of a set according to certain criteria. Specifically, given an equivalence relation on X and x $\in$X, one collects together in the set [x] all elements of X equivalent to x:

[x] := \{y $\in$X : x $\sim$y\}.

The collection of such equivalence classes forms a partition of X (see Sect.2.3.1), given by:

X/$\sim$:= \{[x] : x $\in$X\}. (4.29)

The set X/ $\sim$is called the quotient set of X by $\sim$; it's a set of sets. The idiomatic notation X/$\sim$reminds us of the way this set is constructed: we subdivide X according to the criterion `$\sim$'. The minimalist definition (4.29) is sleek but inefficient, since all the elements of X that belong to the same equivalence class result in just one partition element. So there is a wholesale repetition of partition elements, and the formula works because we agreed to identify repeated elements in set definitions.

The following examples illustrate some applications of this construct.

Example. The relational operator `=' defines an equivalence relation on any set. This is the trivial equivalence; it corresponds to the trivial partition consisting of one-element subsets.

4.6 Relations 71

Example. Given a function f : X $\to$Y, we let x1 $\sim$x2 if f (x1) = f (x2). This is an equivalence relation on X. The equivalence classes are the inverse images of the elements of the co-domain of f :

\{f $-$1(\{y\}) : y $\in$Y\}. (4.30)

Any equivalence relation on a set X can be represented in this way, for some function f .

Example. For any natural number m, the congruence relation x $\equiv$y(mod m) defined in Sect.2.1.3 is an equivalence relation on Z. The equivalence classes are the congruence classes introduced in Sect.2.3.2.

Example. The relation on N $\times$ N, defined by (m, n) $\sim$(j, k) if m + k = n + j, is an equivalence relation. By interpreting the pair (m, n) as the quantity z = m $-$n, we see that equivalent pairs correspond to the same value of z. With this device, one can construct integers from pairs of natural numbers, whereby every integer is an equivalence class of infinitely many pairs of natural numbers. The value of this abstract construction lies with its reductionist character: it requires only natural numbers and addition in N; there is no mention of negative integers or subtraction.

Example. The relation `$\sim$' on Z $\times$ (Z$\setminus$\{0\}), defined by (m, n) $\sim$(j, k) if mk = nj, is an equivalence relation. By interpreting the pair (m, n) as r = m/n, we see that equivalent pairs correspond to the same value of r. With this device, one can define a rational number as an infinite collection of equivalent pairs of integers, without introducing fractions or division.

A relation R on a set X is called a partial ordering if it is reflexive, transitive and anti-symmetric. The latter property is defined as follows:

$\forall$x, y $\in$X, (x R y $\land$y R x) $\Rightarrow$x = y anti-symmetry

A set is partially ordered if a partial ordering is defined on it. A partial ordering is usually denoted by the symbol `$\leq$'. So we write x $\leq$y instead of xRy.

The relational operator $\leq$defines a partial ordering in N, Z, Q, and R (but not in C). The set P(X) of all subsets of a set X is partially ordered by set inclusion, whereby $\leq$means $\subset$.

A partially ordered set X is said to be ordered if all pairs of elements of X are comparable, meaning that we either have x $\leq$y or y $\leq$x. The real line is an ordered set; the power set P(X) of a set X, which is partially ordered by set inclusion, is not ordered.

An ordered set X is said to be well-ordered if any non-empty subset A $\subset$X has a smallest element. The symbolic definition requires three quantifiers:

$\forall$A $\in$P(X)$\setminus$$\emptyset$, $\exists$a $\in$A, $\forall$x $\in$A, a $\leq$x.

Any finite ordered set is well-ordered. The closed unit interval [0, 1] is ordered but not well-ordered, because the subset (0, 1] has no smallest element. The natural

72 4 Mathematical Sentences

numbers are well-ordered. This property forms the basis of the principle of induction, which we consider in Chap.8.

Exercise 4.1 Rewrite each symbolic sentence using the quantifier $\exists$.

1. Y $\supset$X

2. X $\cap$Y = $\emptyset$

3. \#X > 1

4. z $\in$X $\times$ Y

5. X = Y

6. \#(X Y) > 0

7. sin cos = cos sin.

Exercise 4.2 Write each sentence with symbols, using at least one quantifier.

1. The integer n is a cube.

2. The fraction a/b is not reduced.

3. The equation f (x) = 0 has a rational solution.

4. The unit circle has a rational point.

5. The polynomial p(x) has no integer roots.

6. The function f : A $\to$B is not injective.

7. The function f : A $\to$B is not surjective.

8. The function f : A $\to$B is constant.

9. The function f : A $\to$B is not constant.

10. Every cubic equation with real coefficients has a real solution.

11. The integer n has a square divisor.5

12. The integer n is not divisible by 3.

13. Some integers can be written in two different ways as the sum of two cubes.6

14. Any open interval contains a rational number.

15. No open interval contains a smallest element.

16. The equation f (x) = 0 has infinitely many integer solutions.

17. The equation f (x) = 0 has finitely many rational solutions.

Exercise 4.3 Consider the following statements.

1. Every prime greater than 2 is odd.

2. Two integers which are co-prime are also prime.

3. The reciprocal of a positive reduced fraction is reduced.

4. The product of two primes has exactly four divisors.

5. Every divisor of the product of two integers also divides one of the factors.

6. Every subset of an infinite set is infinite.

7. Three consecutive odd integers greater than 3 cannot all be prime.

8. Two complex numbers whose sum is a real number must be complex conjugate of

each other.

9. The inverse of an invertible matrix is invertible.

5 The trivial divisor 12 is excluded. 6 See (4.4).

4.6 Relations 73

In each case:

(i) Decide if the statement is true or false, hence rewrite it as an explicit implication

(For all \dots{}x, if x \dots{}, then \dots{}). (ii) State the contrapositive. (iii) State the converse, and decide whether it's true or false.

(iv) State the negation.

Exercise 4.4 (W. Hodges). We wish to to build up a set of predicates to describe family relations. You are given the two predicates

x is a son of y x is a daughter of y.

Your task is to write definitions of the following predicates, in some appropriate order such that the later definitions use only the given predicates and earlier definitions. (The order below is just alphabetical.)

x is an aunt of y x is a brother of y x is a child of y x is the father of y x is female x is a grandchild of y x is a half-brother of y x is male x is the mother of y x is a nephew of y x is a parent of y x is a sister of y.

Use the symbols x, y, z, etc., to denote people, words for everything else, and parentheses to specify the order of evaluation of logical operators. Why are the given data not quite sufficient to define the `aunt' predicate?

Exercise 4.5 Let x and y be natural numbers, and let P(x, y) mean: `x is a proper divisor of y' (that is, x|y and x = 1, y). Thus P is a predicate over N $\times$ N.

(a) Write each statement with words. [ $\varepsilon$ ]

1. $\exists$x $\in$N, P(x, 5)

2. $\exists$x $\in$N, P(5, x)

3. $\forall$x $\in$N, P(x, x2)

4. $\exists$x $\in$N, $\forall$y $\in$N, P(x, y)

5. $\forall$x $\in$N$\setminus$\{1\}, $\exists$y $\in$N, P(x, y).

74 4 Mathematical Sentences

(b) Write each predicate with words, apart from the predicate's argument. [ $\varepsilon$ ]

1. y $\to$($\exists$x $\in$N, P(x, y))

2. y $\to$($\forall$x $\in$N, $\neg$P(x, y))

3. x $\to$($\exists$y $\in$N, P(x, y)).

(c) Define each set with words. [ $\varepsilon$ ]

1. \{y $\in$N : $\exists$x $\in$N, P(x, y)\}

2. \{y $\in$N : $\forall$x $\in$N, $\neg$P(x, y)\}

3. \{x $\in$N : $\exists$y $\in$N, P(x, y)\}.

Exercise 4.6 The following expressions define sets; turn symbols into words. [ $\varepsilon$ ]

1. \{x $\in$N : $\forall$m, n $\in$N, (x|mn) $\Rightarrow$(x|m $\lor$x|n)\}

2. \{x $\in$N : $\forall$m $\in$N, x|m2 $\Rightarrow$x|m\}

3. \{x $\in$N : $\forall$m $\in$N, x|m3 $\Rightarrow$x|m2\}.

Exercise 4.7 Decide if each sentence is true or false, hence write its negation with symbols. Then rewrite both with words. [ $\varepsilon$ ]

2. $\forall$n $\in$N, $\surd$n $\in$R$\setminus$Q 1. $\forall$n $\in$N, 1/n $\in$N

3. $\forall$x, y $\in$R, xy = yx

4. $\forall$n $\in$Z, 2 | n(n + 1)

5. $\forall$m, n $\in$Z, (m + n $\in$1 + 2Z) $\Rightarrow$(m $\in$2Z $\lor$n $\in$2Z)

6. $\forall$n $\in$N, $\exists$r $\in$Q, n < r2 < n + 1

7. $\forall$y $\in$R, $\exists$x $\in$R+, log(x) = y

8. $\forall$x, y > 1, $\exists$n $\in$N, xn > y

$\surd$

2| < $\epsilon$.

9. $\forall$$\epsilon$ > 0, $\exists$r $\in$Q, |r $-$

Exercise 4.8 The following cryptic symbolic sentence states an interesting mathematical fact:

$\forall$n $\in$N, $\exists$k $\in$N, $\forall$j $\in$N, (k + j)! $\geq$nk+j. (4.31)

Rewrite it in words.

Exercise 4.9 Consider the sentences

1. Everybody loves George

2. Everybody loves somebody

3. Somebody loves everybody.

We wish to determine the computational effort required to find out if a sentence is true or false. With the notation of Sect.4.4, let the set G have n elements. First suppose that the sentence is true. How many evaluations of the predicate `x loves y' are needed to verify this? Give the minimum and maximum number of evaluations, with an explanation. Do the same assuming that the sentence is false.

4.6 Relations 75

Exercise 4.10 In this exercise we collect some straightforward proofs.

1. Show that the operators $\lor$, $\Rightarrow$, $\Leftrightarrow$may be defined in terms of $\neg$ and $\land$.

2. Complete the proof of Theorems 4.1, 4.2, and 4.3.

3. Show that given any predicates P, Q over a set X, there is a subset A of X such

that

($\forall$x $\in$X, P(x) $\Rightarrow$Q(x)) $\Leftrightarrow$($\forall$x $\in$A, Q(x)).

This equivalence says that the hypothesis of an implication may be absorbed into a suitable ambient set.

Exercise 4.11 Define some interesting predicates on the power set of Z. Do the same for the power set of Z[x].

Exercise 4.12 Let X be any set. Show that there is a one-to-one correspondence between the partitions of X and the equivalence relations on X. Show that any equivalence relation on X may be represented by a suitable function f : X $\to$Y such that the equivalence classes are the pre-images of the elements of Y---see Eq. (4.30).

Exercise 4.13 You are given an equivalence relation over a finite set X. Develop an algorithm---more efficient than formula (4.29)---to construct the associated partition of X.

\end{document}